\documentclass[a4paper]{article}

\usepackage{latexml}
\usepackage[T1, T5]{fontenc}
\usepackage[utf8]{inputenc}
\usepackage{amsmath}
\usepackage{amsthm}
\usepackage{amssymb}
\usepackage{tikz}
\usetikzlibrary{shapes}
\usepackage[american,vietnamese]{babel}
\usepackage{algorithm, algpseudocode}
\usepackage{graphicx}
\usepackage{tabularx}
\newcolumntype{Y}{>{\centering\arraybackslash}X} % center align in X column (using Y instead of X)
\usepackage{multirow}
\usepackage{array}
\usepackage{listings}
\usepackage{color}
\usepackage[margin=1in]{geometry}
\usepackage{newtxtext,newtxmath}

\usepackage[sort&compress,numbers,square,comma]{natbib}
%\usepackage[%
%			backend=biber, 
%			bibstyle=numeric, 
%			citestyle=numeric-comp,
%			maxcitenames=3,
%			maxbibnames=20,
%			sorting=ydnt
%			]{biblatex}
%\addbibresource{refs}

\definecolor{mygreen}{rgb}{0,0.6,0}
\definecolor{mygray}{rgb}{0.5,0.5,0.5}
\definecolor{mymauve}{rgb}{0.58,0,0.82}
\definecolor{mybackground}{rgb}{0.83,0.89,0.88}
\definecolor{darkgray}{rgb}{0.66, 0.66, 0.66}
\definecolor{x11gray}{rgb}{0.75, 0.75, 0.75}
\definecolor{internallinkcolor}{rgb}{0,.5,0}
\definecolor{externallinkcolor}{rgb}{0,0,.5}

\usepackage[unicode,
			colorlinks,
			linktoc=all, % defines which part of an entry in the table of contents is made into a link  
			linkcolor=externallinkcolor,	% color of internal links (sections, pages, etc.), default is red
			citecolor=internallinkcolor, 	% color of citation links (bibliography), default is green
			filecolor=externallinkcolor, % % color of file links, default is cyan
			urlcolor=externallinkcolor   % color of URL links (mail, web), default is magneta
			]{hyperref}

\lstset{ 
  backgroundcolor=\color{white},   % choose the background color; you must add \usepackage{color} or \usepackage{xcolor}; should come as last argument
  basicstyle=\ttfamily\small,        % the size of the fonts that are used for the code
  breakatwhitespace=false,         % sets if automatic breaks should only happen at whitespace
  breaklines=true,                 % sets automatic line breaking
  captionpos=b,                    % sets the caption-position to bottom
  commentstyle=\color{mygray},    % comment style
  deletekeywords={...},            % if you want to delete keywords from the given language
  escapeinside={\%*}{*)},          % if you want to add LaTeX within your code
  extendedchars=true,              % lets you use non-ASCII characters; for 8-bits encodings only, does not work with UTF-8
  frame=single,	                   % adds a frame around the code
  keepspaces=true,                 % keeps spaces in text, useful for keeping indentation of code (possibly needs columns=flexible)
  keywordstyle=\color{blue},       % keyword style
%  language=Octave,                 % the language of the code
  morekeywords={*,...},            % if you want to add more keywords to the set
  numbers=left,                    % where to put the line-numbers; possible values are (none, left, right)
  numbersep=5pt,                   % how far the line-numbers are from the code
  numberstyle=\small\color{black}, % the style that is used for the line-numbers
  rulecolor=\color{black},         % if not set, the frame-color may be changed on line-breaks within not-black text (e.g. comments (green here))
  showspaces=false,                % show spaces everywhere adding particular underscores; it overrides 'showstringspaces'
  showstringspaces=false,          % underline spaces within strings only
  showtabs=false,                  % show tabs within strings adding particular underscores
  stepnumber=1,                    % the step between two line-numbers. If it's 1, each line will be numbered
  stringstyle=\color{mymauve},     % string literal style
  tabsize=2,	                   % sets default tabsize to 2 spaces
  title=\lstname                   % show the filename of files included with \lstinputlisting; also try caption instead of title
}

\usepackage{xcolor}

% Solarized colors
%\definecolor{sbase03}{HTML}{002B36}
%\definecolor{sbase02}{HTML}{073642}
%\definecolor{sbase01}{HTML}{586E75}
%\definecolor{sbase00}{HTML}{657B83}
%\definecolor{sbase0}{HTML}{839496}
%\definecolor{sbase1}{HTML}{93A1A1}
%\definecolor{sbase2}{HTML}{EEE8D5}
%\definecolor{sbase3}{HTML}{FDF6E3}
%\definecolor{syellow}{HTML}{B58900}
%\definecolor{sorange}{HTML}{CB4B16}
%\definecolor{sred}{HTML}{DC322F}
%\definecolor{smagenta}{HTML}{D33682}
%\definecolor{sviolet}{HTML}{6C71C4}
%\definecolor{sblue}{HTML}{268BD2}
%\definecolor{scyan}{HTML}{2AA198}
%\definecolor{sgreen}{HTML}{859900}
%
%\lstset{ % Solarized-Dark
%    % How/what to match
%    sensitive=true,
%    % Border (above and below)
%    frame=lines,
%    % Extra margin on line (align with paragraph)
%    xleftmargin=\parindent,
%    % Put extra space under caption
%    belowcaptionskip=1\baselineskip,
%    % Colors
%    backgroundcolor=\color{sbase03},
%    basicstyle=\color{sbase00}\ttfamily,
%    keywordstyle=\color{scyan},
%    commentstyle=\color{sbase1},
%    stringstyle=\color{sblue},
%    numberstyle=\color{sviolet},
%    identifierstyle=\color{sbase00},
%    %identifierstyle=
%    % Break long lines into multiple lines?
%    breaklines=true,
%    % Show a character for spaces?
%    showstringspaces=false,
%    tabsize=2
%}

%\lstset{ % Solarized-Light
%    % How/what to match
%    sensitive=true,
%    % Border (above and below)
%    frame=lines,
%    % Extra margin on line (align with paragraph)
%    xleftmargin=\parindent,
%    % Put extra space under caption
%    belowcaptionskip=1\baselineskip,
%    % Colors
%    backgroundcolor=\color{sbase3},
%    basicstyle=\color{sbase00}\ttfamily,
%    keywordstyle=\color{scyan},
%    commentstyle=\color{sbase1},
%    stringstyle=\color{sblue},
%    numberstyle=\color{sviolet},
%    identifierstyle=\color{sbase00},
%    % Break long lines into multiple lines?
%    breaklines=true,
%    % Show a character for spaces?
%    showstringspaces=false,
%    tabsize=2
%}

\newtheorem{theorem}{Theorem}
\newtheorem{lemma}[theorem]{Lemma}
\newtheorem{conjecture}[theorem]{Conjecture}
\newtheorem{problem}[theorem]{Open Problem}

\theoremstyle{definition}
\newtheorem{definition}[theorem]{Definition}
\newtheorem{example}[theorem]{Example}
\newtheorem{exercise}[theorem]{Exercise}

\theoremstyle{remark}
\newtheorem{remark}[theorem]{Remark}

\numberwithin{equation}{section}

\hypersetup{
	pdfauthor={Duc A. Hoang},
	pdftitle={Open Problems},
	pdfsubject={Open Problems},
	pdfkeywords={open problem, graph algorithm, reconfiguration problem},
	pdfproducer={LaTeX},
	pdfcreator={pdfLaTeX}
}

\iflatexml
% no title and author in Jekyll posts
\title{}
\author{}
\date{\today}
\else
\title{\textbf{Open Problems}}
\author{Duc A. Hoang\\ VNU University of Science, Hanoi, Vietnam\\ \href{mailto:hoanganhduc@hus.edu.vn}{\texttt{hoanganhduc@hus.edu.vn}}}
\fi

\begin{document}

\nocite{*}

\selectlanguage{american}

\maketitle

\begin{abstract}
This document contains a list of open problems that \href{https://hoanganhduc.github.io/}{Duc A. Hoang} is interested in.  
It is also available in
%
\iflatexml 
\href{https://hoanganhduc.github.io/problems/main.pdf}{PDF} 
\else 
\href{https://hoanganhduc.github.io/problems/}{HTML} 
\fi%
format.
Please get in touch if you have any idea/comment/question/update.
%
You can find (probably more interesting) open problems from the following resources.
\begin{itemize}
	\item \href{http://www.openproblemgarden.org/category/graph_theory}{Open Problem Garden: Graph Theory}.
	\item \href{https://faculty.math.illinois.edu/~west/openp/}{Open Problems - Graph Theory and Combinatorics}, by \href{https://math.uiuc.edu/~west/}{Douglas B. West}.
	\item \href{https://nms.kcl.ac.uk/iwoca/problems_list.html}{List of IWOCA's Open Problems}.
	\item \href{https://link.springer.com/book/10.1007/978-3-319-31940-7}{Graph Theory - Favorite Conjectures and Open Problems - 1} (edited by Ralucca Gera, Stephen Hedetniemi, Craig Larson) and \href{https://link.springer.com/book/10.1007/978-3-319-97686-0}{Graph Theory - Favorite Conjectures and Open Problems - 2} (edited by Ralucca Gera, Teresa W. Haynes, Stephen T. Hedetniemi).
	\item \href{https://topp.openproblem.net}{Open Problems in Computational Geometry}, maintained by \href{http://erikdemaine.org/}{Erik D. Demaine}, \href{http://www.ams.sunysb.edu/~jsbm/}{Joseph S. B. Mitchell}, \href{http://cs.smith.edu/~orourke/}{Joseph O’Rourke}.
	\item \href{http://fpt.wikidot.com/open-problems}{Open Problems in Paramaterized Complexity}.
	\item \href{https://en.wikipedia.org/wiki/Category:Unsolved_problems_in_graph_theory}{List of unsolved problems/conjectures in graph theory from Wikipedia}.
\end{itemize}

\end{abstract}

\tableofcontents

\section{Preliminaries}

For the terminology and notation not defined here, see the textbooks by Diestel~\cite{Diestel2017} or West~\cite{West2001}. 
We often denote by $P_n, C_n, K_n, K_{m,n}$ the \textit{path}, \textit{cycle}, \textit{complete graph}, \textit{complete bipartite graph} on $n$ vertices, respectively.
For two graphs $G, H$, the \textit{disjoint union} of $G$ and $H$, denoted by $G + H$, is the graph with $V(G+H) = V(G) \cup V(H)$ and $E(G+H) = E(G) \cup E(H)$.

\section{Reconfiguration of Independent Sets under Token Sliding}

\subsection{Definitions}

Let $G$ be a simple, undirected graph.
The \textit{double-broom graph} $D_{r,n,s}$ is the tree obtained from a path $P_n = v_1v_2\dots v_n$ by attaching $r$ leaves to $v_1$ and $s$ leaves to $v_n$, where $r, n, s$ are positive integers. 
Two independent sets $I, J$ are \textit{adjacent under Token Sliding $(\mathsf{TS})$} if there exists $u, v \in V(G)$ such that $I \setminus J = \{u\}$, $J \setminus I = \{v\}$, and $uv \in E(G)$.
The \textsc{Token Sliding} problem asks whether there is a \textit{$\mathsf{TS}$-sequence} between $I$ and $J$ in $G$, that is, a sequence of independent sets starting with $I$ and ending with $J$ where any two consecutive members are adjacent under $\mathsf{TS}$.
For a fixed positive integer $k$, the graph $\mathsf{TS}_k(G)$ contains all size-$k$ independent sets as its vertices/nodes, and two nodes are joined by an edge if their corresponding independent sets are adjacent under $\mathsf{TS}$.

\subsection{Conjecture/Question}

\begin{problem}\label{ISR-TS-tw2graph}
	What is the complexity of \textsc{Token Sliding} on outerplanar graphs? And more generally, on series-parallel graphs ($\equiv$ graphs of treewidth at most two)?
\end{problem}

\begin{conjecture}\label{ISR-TSkgraph}
	Let $G$ be a forest.
	Then $\mathsf{TS}_k(G)$ is a forest if and only if $G$ is $\{2K_2 + (k-2)K_1, D_{2,2,2} + (k-2)K_1, D_{2,4,2} + (k-3)K_1\}$-free, for some integer $k \geq 4$.
\end{conjecture}

\subsection{Background/Motivation}

\subsection{Comments/Partial Results}

\section{Reconfiguring $k$-Path Vertex Covers on Trees under Token Sliding}

\bibliographystyle{mplainnat}
\bibliography{refs}
%\printbibliography

\end{document}
